\section*{第 \ref{ch:g11:antibiotic-resistance} 章——细菌与抗生素耐药性}

\begin{solution}{exo:g11:antibiotic-resistance:1}
克隆分裂过程中产生的\omterm{def:g11:mutations:mutation}{突变}，以及
水平\omterm{def:g10:universal-dna:gene}{基因}转移——即从另一细菌获得的 \omterm{def:g10:universal-dna:information}{DNA}，通常是质粒。
\end{solution}

\begin{solution}{exo:g11:antibiotic-resistance:2}
一类分子，通过作用于细菌具有而人体\omterm{def:g10:cells-common-unit:cell}{细胞}没有的结构，
杀死细菌或抑制其生长（如\omterm{def:g10:cells-common-unit:organelle}{细胞壁}、细菌\omterm{def:g10:cells-common-unit:organelle}{核糖体}、
细菌复制酶）。感冒由病毒引起，
而病毒没有上述任何靶点。
\end{solution}

\begin{solution}{exo:g11:antibiotic-resistance:3}
一种能够分解\omterm{def:g11:antibiotic-resistance:antibiotic}{抗生素}的酶；一个发生\omterm{def:g11:mutations:mutation}{突变}、药物不再能与之结合的靶点；
一种将药物排出的泵（或能够阻挡药物进入的\omterm{def:g10:cells-common-unit:organelle}{细胞壁}）。
\end{solution}

\begin{solution}{exo:g11:antibiotic-resistance:4}
敏感（20 高于 18），但仅略高于阈值。
D 高出阈值的幅度最大，其次是 A。
\end{solution}

\begin{solution}{exo:g11:antibiotic-resistance:5}
否。在影印平板上，耐药菌落在每一块含\omterm{def:g11:antibiotic-resistance:antibiotic}{抗生素}的平板上都出现在相同位置，
因此耐药\omterm{def:g10:cells-common-unit:cell}{细胞}早已存在于原始菌落中，
而这些菌落从未接触过药物。\omterm{def:g11:antibiotic-resistance:antibiotic}{抗生素}
只是显现并选择耐药\omterm{def:g10:cells-common-unit:cell}{细胞}，而不是诱导其产生。
\end{solution}

\begin{solution}{exo:g11:antibiotic-resistance:6}
30 次分裂：$2^{30} \approx 10^9$ 个\omterm{def:g10:cells-common-unit:cell}{细胞}。
约 $10^9$ 次分裂产生了它们：平均约 10 个耐药\omterm{def:g10:cells-common-unit:cell}{细胞}——
若\omterm{def:g11:mutations:mutation}{突变}出现得较早，则数目更多。
\end{solution}

\begin{solution}{exo:g11:antibiotic-resistance:7}
约第 3 天，此时敏感\omterm{def:g10:cells-common-unit:cell}{细胞}降至约 0.2\%，耐药\omterm{def:g10:cells-common-unit:cell}{细胞}升至约 1\%。
第 6 天时，虚线所示的总量恢复到 100\%，
几乎全部由耐药\omterm{def:g10:cells-common-unit:cell}{细胞}构成（第 3 天存活的敏感\omterm{def:g10:cells-common-unit:cell}{细胞}也重新生长，
但起点较小）。
\end{solution}

\begin{solution}{exo:g11:antibiotic-resistance:8}
位于\omterm{def:g11:cell-cycle-mitosis:chromosome}{染色体}上的\omterm{def:g10:universal-dna:gene}{基因}只传递给克隆的后代；
质粒还可以通过水平转移传递给无亲缘关系的\omterm{def:g10:cells-common-unit:cell}{细胞}和其他\omterm{def:g10:biodiversity-scales:species}{物种}，
并且一次转移就能使\omterm{def:g10:cells-common-unit:cell}{细胞}立即获得耐药性，
无需等待\omterm{def:g11:mutations:mutation}{突变}。
\end{solution}

\begin{solution}{exo:g11:antibiotic-resistance:9}
第 3 天时，敏感\omterm{def:g10:cells-common-unit:cell}{细胞}已降至百分之几以下，
而耐药\omterm{def:g10:cells-common-unit:cell}{细胞}正在增加；若停药，所有存活者
都会重新生长，且恢复的种群中耐药\omterm{def:g10:cells-common-unit:cell}{细胞}的
比例远高于原始种群。复发的感染
不再对同一种药物有效。
\end{solution}

\begin{solution}{exo:g11:antibiotic-resistance:10}
医院对许多相邻的患者密集使用\omterm{def:g11:antibiotic-resistance:antibiotic}{抗生素}：选择压力持续存在，
被选出的菌株
借助双手和物体表面在患者之间传播。
\end{solution}

\begin{solution}{exo:g11:antibiotic-resistance:11}
链霉素通过与\omterm{def:g10:cells-common-unit:organelle}{核糖体}蛋白结合发挥作用；
\omterm{def:g11:gene-expression:protein}{蛋白质}发生改变后不再与之结合，\omterm{prop:g11:gene-expression:translation}{翻译}得以继续。
但这种改变也会轻微损害\omterm{def:g10:cells-common-unit:organelle}{核糖体}的正常功能，
因此\omterm{def:g11:gene-expression:protein}{蛋白质}合成及生长都会略慢——这是耐药性的一种代价。
\end{solution}

\begin{solution}{exo:g11:antibiotic-resistance:12}
每次分裂的概率为 $10^{-8} \times 10^{-8} = 10^{-16}$。在 $10^{12}$
个细菌中，几乎没有\omterm{def:g10:cells-common-unit:cell}{细胞}可能同时携带两者。结核病的治疗
持续数月，细菌数量可达数十亿，因此单一药物的耐药性
必然被选出；同时使用两种药物，则不留下可供选择的
存活者。
\end{solution}

\begin{solution}{exo:g11:antibiotic-resistance:13}
低剂量会在动物肠道中选择出耐药菌株，而不会杀死整个种群；
工人通过接触获得这些细菌；质粒将\omterm{def:g10:universal-dna:gene}{基因}
传递给人类病原体，并经粪肥进入环境。
对该药物及相关的人用药物的耐药性
在医学上随之上升。
\end{solution}

\begin{solution}{exo:g11:antibiotic-resistance:14}
耐药性往往让细菌付出代价（如\omterm{def:g10:cells-common-unit:organelle}{核糖体}速度变慢、需要合成某种酶）：
在无药物条件下，敏感\omterm{def:g10:cells-common-unit:cell}{细胞}生长稍快，
经过多代便能在竞争中胜过耐药\omterm{def:g10:cells-common-unit:cell}{细胞}。耐药
\omterm{def:g10:cells-common-unit:cell}{细胞}的比例随之下降，但耐药\omterm{def:g10:universal-dna:gene}{基因}仍以低频率
留存于种群中，随时可能再次被选择。
\end{solution}

\begin{solution}{exo:g11:antibiotic-resistance:15}
细菌并非个体适应。在任何大种群中，
都有少数\omterm{def:g10:cells-common-unit:cell}{细胞}，通过随机\omterm{def:g11:mutations:mutation}{突变}或从另一\omterm{def:g10:cells-common-unit:cell}{细胞}获得的质粒，
携带使其产生耐药性的\omterm{def:g10:universal-dna:gene}{等位基因}。\omterm{def:g11:antibiotic-resistance:antibiotic}{抗生素}杀死
其他\omterm{def:g10:cells-common-unit:cell}{细胞}；耐药\omterm{def:g10:cells-common-unit:cell}{细胞}增殖，并将该\omterm{def:g10:universal-dna:gene}{等位基因}
传递给后代。种群之所以变得耐药，是因为它被
进行了筛选，而非其成员学会了什么。
\end{solution}

\begin{solution}{pb:g11:antibiotic-resistance:1}
\textbf{1.} 约 $10^{10}$ 次分裂（每产生一个\omterm{def:g10:cells-common-unit:cell}{细胞}对应一次）。

\textbf{2.} 平均而言，$10^{10} \times 10^{-8} = 100$ 个耐药\omterm{def:g10:cells-common-unit:cell}{细胞}。

\textbf{3.} 一百亿个\omterm{def:g10:cells-common-unit:cell}{细胞}中的一百个在平板上不可见：
菌苔之所以表现为敏感，是因为其中 99.999999\% 为敏感\omterm{def:g10:cells-common-unit:cell}{细胞}。

\textbf{4.} 阿莫西林：敏感（24 大于 17）。红霉素：
耐药（12 小于 20）。青霉素：完全耐药，无抑菌作用。

\textbf{5.} 青霉素对该菌株无效。在敏感药物中选择时，
还需权衡副作用与抗菌谱；
应优先选用有效且抗菌谱最窄的药物，以保护患者体内的其他
细菌。

\textbf{6.} 每天四个周期，每个周期 6 小时：$10^{10} \times 0.1^4 =
10^6$，1 天后；$10^{10} \times 0.1^{12} = 0.01$，3 天后——
基本为零。

\textbf{7.} $100 \times 2^{48} \approx \num{2.8e16}$：荒谬，超过
伤口所能容纳的数量。生长受空间、营养物质和
免疫系统的限制。

\textbf{8.} 从 100 个\omterm{def:g10:cells-common-unit:cell}{细胞}增至 $10^{10}$ 个，需要 $\log_2 10^8 \approx 27$
次倍增，约 13 小时；实际上，空间在第一天内逐渐腾出，
因此到第 1 至 2 天时，耐药克隆已取代敏感
种群。

\textbf{9.} 敏感菌：每 6 小时数量降为原来的十分之一，直至
第 3 天降至零。耐药菌：从 100 个增长，
一两天内占满空间，随后稳定在 $10^{10}$。总数：第 1 天有所下降，之后
回升至 $10^{10}$，且全部为耐药菌。

\textbf{10.} 若免疫系统清除\omterm{def:g10:cells-common-unit:cell}{细胞}的速度快于
这一百个耐药\omterm{def:g10:cells-common-unit:cell}{细胞}增殖的速度，感染仍会终止；若
克隆的增长超过清除速度，感染便转为耐药。
\omterm{def:g11:antibiotic-resistance:antibiotic}{抗生素}是与免疫系统共同发挥作用，而非取代免疫系统。

\textbf{11.} 2 天后的敏感菌：$10^{10} \times 0.1^8 = 100$。
耐药菌：若克隆已占满空间，约为 $10^{10}$；否则
至少为数千个。存活者几乎全部为耐药菌。

\textbf{12.} 以相同速率增长，比例保持不变：
基本为 100\% 耐药。

\textbf{13.} 无效：耐药克隆不受阿莫西林影响。

\textbf{14.} 第一个平板上是原始种群，其中 99.999999\%
为敏感菌；复发平板上是被选择出的克隆，全部为耐药菌，
因此菌苔一直生长至药敏纸片处。

\textbf{15.} 该\omterm{def:g10:universal-dna:gene}{基因}可传递给由其他患者或工作人员
携带的其他\omterm{def:g10:biodiversity-scales:species}{物种}，因此病房中的其他病原体
无需等待自身\omterm{def:g11:mutations:mutation}{突变}即可产生耐药性。

\textbf{16.} 在病房中，每年 400 个疗程作用于一个小型封闭
种群，持续对细菌进行筛选，被选中的菌株
在患者之间传播；在城镇中，选择压力分散于众多
人群与\omterm{def:g10:biodiversity-scales:species}{物种}之间。

\textbf{17.} 药敏试验：选用菌株尚无耐药性的药物，
从而不会筛选出克隆。完整疗程：不中断疗程，以免
耐药存活者再度增殖。卫生：阻断被选中克隆在患者间
的传播。隔离：阻断携带者造成的水平或
垂直传播。

\textbf{18.} 选择压力减小后，生长稍快的敏感菌株
得以恢复优势，耐药菌株的传播也被切断。其
比例未回到 5\%，是因为耐药菌株及其质粒
仍以低频率存在，并被每个
剩余疗程重新选中。

\textbf{19.} 任何大型细菌种群中，都有因偶然\omterm{def:g11:mutations:mutation}{突变}而
对新药耐药的\omterm{def:g10:cells-common-unit:cell}{细胞}；使用该药即会将其筛选出来，数年之内
耐药菌株便会普遍存在，此前引入的每种\omterm{def:g11:antibiotic-resistance:antibiotic}{抗生素}
均是如此。

\textbf{20.} 在用药之前，一百亿个\omterm{def:g10:cells-common-unit:cell}{细胞}中约有一百个
耐药\omterm{def:g10:cells-common-unit:cell}{细胞}；治疗第二天，它们已占据伤口；
过早停止疗程，使本可治愈的感染转变为
扩散的耐药感染。
\end{solution}
